\section{Introduction}
Image denoising, the recovery of a clean image from a corrupted observation, is a basic problem in computer vision and the core learning step of denoising diffusion models~\cite{ho}. Because operating directly on pixels is computationally expensive, latent-space methods first compress images with an autoencoder and then process the compact representation~\cite{rombach}. In parallel, variational quantum algorithms have made parameterized quantum circuits (PQCs) a common building block of hybrid quantum--classical learning~\cite{cerezo,benedetti}, raising the question of whether such circuits can serve as useful components of image reconstruction models.

The specific problem addressed in this work is the removal of Gaussian noise from the latent representation of an image using a model that contains a PQC. Noisy intermediate-scale quantum (NISQ) devices provide few qubits and are sensitive to noise~\cite{preskill}, whereas even the latent map of a $28\times28$ image in our autoencoder holds $64\times7\times7=3{,}136$ values. The representation must therefore be divided and strongly compressed before it can be encoded into a circuit. This compression, the data encoding and the measurement scheme all limit what the quantum layer can express~\cite{havlicek,schuld}, while deeper circuits can be hard to train~\cite{mcclean}.

Quantum circuits have already been incorporated into generative and denoising models. Quantum diffusion models have been studied for image data~\cite{kolle,cacioppo} and quantum data~\cite{zhang}, and patch-based designs have enabled image generation with few qubits~\cite{huang}. Showing that a hybrid model works, however, does not reveal which of its parts is responsible for the result. Here, the classical layers contain thousands of times more trainable parameters than the circuit. A large benchmarking study of quantum classifiers further found that conclusions depend strongly on the choice of baselines and that removing quantum properties does not necessarily reduce performance~\cite{bowles}. For patch-based latent denoising, the effects of circuit depth, encoding and readout, and the dependence of the model on its quantum layer, have not been examined closely.

This gap matters because quantum resources are scarce and costly to simulate or access. Without strong baselines and ablation tests, an improvement caused by a wider classical network, a longer training schedule or an extra refinement stage may wrongly be credited to the circuit. Careful attribution helps researchers decide which design choices deserve further investigation.

We therefore present PHQD, a patch-based hybrid quantum--classical framework for latent image denoising on Fashion-MNIST~\cite{fmnist}. A convolutional autoencoder produces the latent map, which is corrupted with Gaussian noise and divided into overlapping patches. A classical network converts each patch into rotation angles, a PQC processes them, and the output is mapped back to patch space; the patches are then recombined and decoded.

The objectives of this study are:
\begin{enumerate}
\item To design a modular hybrid pipeline for patch-wise denoising of latent image representations.
\item To compare the effects of circuit depth, data encoding and measurement strategy across several quantum configurations.
\item To determine, using classical baselines and ablation tests, how strongly performance depends on the quantum layer.
\end{enumerate}

The main contributions of this paper are:
\begin{itemize}
\item An end-to-end implementation combining a convolutional autoencoder, overlapping patch processing, a batched PQC denoiser and a stitching network.
\item A comparison of four 4-qubit configurations with a 6-qubit model that uses data re-uploading, three-basis measurement and joint refinement.
\item An evaluation on the full test set against multilayer perceptron, U-Net and linear regression baselines.
\item An ablation analysis showing that reliance on the circuit output differs considerably between configurations, together with a stitching-network control indicating that part of the 6-qubit model's advantage may arise from classical refinement.
\end{itemize}

The rest of this paper is organized as follows. Section~\ref{sec:related} reviews related work. Section~\ref{sec:method} describes the PHQD architecture and the quantum configurations. Section~\ref{sec:setup} presents the experimental setup and baselines. Section~\ref{sec:results} reports the results and ablation analysis, and Section~\ref{sec:conclusion} concludes the paper.
